\chapter{La toile sur la rivière}

\sessioninfo{Session 8}{Mardi 6 octobre 2026}

Le canoë occupait toujours le fond du couloir de la maison de repos de Kaya. Les quatre aventuriers avaient choisi leur guide, acheté leur matériel et fixé leur départ au lendemain matin. L'après-midi touchait à sa fin, mais une dernière question continuait de les préoccuper : étaient-ils réellement prêts à affronter les poisons, les maladies et les caprices de la jungle ?

Ils décidèrent de retourner une fois encore à la capitainerie. Dans les rues de Port Nyanzaru, les passants commençaient déjà à regagner leurs demeures et les étals fermaient progressivement. Morrÿs fit remarquer que le groupe passait beaucoup de temps à parcourir la ville dans tous les sens.

\begin{narrateur}

Quand on n'a pas de tête, il faut avoir des jambes.

\end{narrateur}

Le Sang-dragon doré rencontré plus tôt les reconnut immédiatement. Il approuva leur choix lorsqu'ils lui annoncèrent avoir engagé Azaka : à ses yeux, elle connaissait suffisamment la jungle pour leur offrir de bonnes chances d'en revenir. Obélik l'interrogea alors sur la météo, tandis que Kael revint sur les poisons et les maladies susceptibles de les menacer.

Prévoir précisément le temps au-delà de la côte relevait presque de l'impossible. Les pluies étaient fréquentes, parfois plusieurs fois par jour, mais généralement courtes. Les véritables dangers venaient des orages plus violents : les rivières et les ruisseaux pouvaient enfler brusquement, emportant tout sur leur passage dans des courants impossibles à maîtriser.

Le groupe échangea un regard.

Ils venaient précisément d'acheter un canoë.

Pour les maladies, le Sang-dragon ne possédait aucun remède universel à leur conseiller. La meilleure protection consistait encore à ne rien contracter, conseil irréprochable dont l'utilité pratique restait plus discutable. Il leur recommanda surtout de faire confiance à leur guide, de rester prudents et de ne jamais sous-estimer la jungle.

Les morts-vivants constituaient un autre risque. Ils pouvaient apparaître presque partout, mais devenaient plus nombreux à mesure que l'on pénétrait sur leurs territoires. Leur lenteur et leur incapacité à nager offraient toutefois un avantage appréciable à un groupe voyageant par voie d'eau.

Avant de les laisser repartir, le Sang-dragon leur parla d'une rumeur circulant sur les quais. Un dirigeable nommé le \textit{Star Goddess} aurait dû atteindre Port Nyanzaru depuis quelque temps, mais personne ne l'avait vu arriver. S'il avait atteint Chult, il s'était peut-être écrasé quelque part sur l'île. Dans le cas contraire, son sort pouvait être bien pire encore.

Les aventuriers promirent de garder les yeux ouverts. Leur interlocuteur leur souhaita de revenir en bonne santé, riches et couverts de succès. Kael le remercia avec sincérité : le revoir signifierait au moins qu'ils seraient effectivement revenus.

\scenebreak

La maison de repos était bien plus animée lorsqu'ils y retournèrent. Kaya courait d'une table à l'autre, chargée de plats et de boissons, tandis que les conversations remplissaient la salle. Le groupe trouva une place près de l'entrée, les tables proches du feu étant déjà occupées.

Obélik chercha Azaka du regard. Leur guide n'était pas là, mais une femme assise seule au fond de la salle observait le groupe avec une attention difficile à ignorer. Brune, la peau mate et les cheveux attachés en une longue queue-de-cheval, elle portait autour d'un œil un tatouage qui durcissait encore son regard sévère.

Kael finit par la reconnaître. Elle figurait sur l'un des prospectus consacrés aux guides de Port Nyanzaru. Il s'agissait de Salida.

\illustration[0.9\textwidth]{2026-10-06/0d1aad53-12c9-41d4-a560-72d18abb52b5.png}
{\textit{Salida, guide de Port Nyanzaru.}}
{salida}

Kaya vint prendre leur commande avant qu'ils n'aient décidé comment réagir. Le repas couvert par Wakanga se composait ce soir-là d'une soupe de langoustes au curry vert, préparée avec du lait de coco, de l'ail et suffisamment de piment pour rappeler aux voyageurs qu'ils avaient quitté la Côte des Épées. La description fut assez détaillée pour que certains n'en comprennent pas chaque ingrédient.

Cela avait tout de même l'air excellent.

Pour accompagner le repas, Kaya leur proposa du tej, une boisson locale fermentée. Obélik vida presque immédiatement sa chope et découvrit une saveur sucrée et florale, proche d'un hydromel légèrement pétillant. Destos préféra prendre son temps. Kael, qui n'appréciait guère l'alcool, accepta malgré tout d'y goûter.

Kaya revint bientôt avec les soupes et une assiette supplémentaire. Les nachos étaient offerts par la femme installée au fond de la salle et accompagnés d'un petit mot. Salida y vantait ses compétences et se disait certaine que la guide choisie par le groupe ne pourrait rivaliser avec elle. Elle leur demandait s'ils accepteraient au moins de discuter.

Destos refusa de toucher à l'assiette. Obélik se proposa aussitôt de manger sa part.

Kael salua Salida de la main pour la remercier. La guide interpréta le geste comme une invitation, rejoignit leur table et demanda poliment la permission de s'asseoir. Elle se présenta comme une professionnelle parcourant Chult depuis plus de trente ans et affirma que presque tous les groupes qu'elle avait accompagnés étaient revenus vivants, généralement chargés de richesses.

Elle ne manquait pas de confiance. Salida se considérait comme la plus polyvalente et la plus compétente des guides disponibles. Contrairement à ceux qui entraînaient leurs clients dans leurs propres affaires, elle promettait de se consacrer entièrement à la mission pour laquelle on l'engageait.

\illustration[0.9\textwidth]{2026-10-06/171929ed-6cfb-4bb1-a62c-f959d53c1316.png}
{\textit{Salida tenta de convaincre les aventuriers de faire appel à ses services.}}
{rencontre_salida}

Le groupe lui apprit qu'il avait déjà donné sa parole à Azaka. Ses ressources limitées avaient fortement pesé dans ce choix : les honoraires réglementaires, payables trente jours à l'avance, dépassaient même l'or qu'il leur restait. Salida expliqua qu'elle ne pouvait ni réduire son tarif ni offrir ses services sans créer une concurrence déloyale envers les autres guides.

Destos observa attentivement son attitude. Malgré son assurance et ses critiques à peine voilées, elle semblait sincèrement convaincue de ses compétences et de son professionnalisme. Elle travaillait pour l'argent, bien sûr, mais affirma que la mission de ses clients passerait toujours avant ses intérêts personnels.

Les aventuriers demandèrent quelques instants pour en discuter. Ils s'étaient déjà engagés auprès d'Azaka et refusaient de revenir sur leur parole à la veille du départ.

Destos alla lui annoncer leur décision. Salida l'accepta sans insister davantage, regrettant seulement qu'ils accordent un prix limité à leur propre sécurité. Elle leur souhaita de prendre soin d'eux, puis quitta la maison de repos.

Le groupe put enfin terminer son repas.

La soirée s'acheva paisiblement. Les clients regagnèrent peu à peu leurs chambres ou quittèrent l'établissement, tandis que Kaya nettoyait les dernières tables. Elle se réjouit d'apprendre que le canoë disparaîtrait de son couloir dès le lendemain matin.

\scenebreak

La nuit fut calme pour presque tout le monde. Le vent soufflait contre les murs et les ronflements des voyageurs traversaient les cloisons, sans que rien ne vienne troubler davantage l'auberge. Kael, lui, dormit mal. Dans son demi-sommeil, chaque rafale prenait l'ampleur d'un ouragan et chaque craquement annonçait une vague prête à engloutir le bâtiment.

Le jour se leva pourtant sur un temps agréable. Une humidité constante chargeait l'air, mais la température restait douce et une légère brise descendait du nord. Kael frappa à la porte de Morrÿs et d'Obélik pendant que Destos sortait vérifier si Azaka les attendait déjà.

Elle était adossée au panneau des guides, jouant distraitement avec son collier de perles. À peine aperçut-elle Destos qu'elle s'étonna de leur manque d'empressement. Le clerc lui proposa de les rejoindre pour le petit déjeuner, mais Azaka avait déjà mangé et leur rappela qu'une longue route les attendait.

Le repas du matin fut abandonné. Morrÿs rappela au groupe de ne pas oublier le canoë, tandis que Destos rendait les clés des chambres à Kaya. La tenancière leur souhaita une bonne exploration et espéra les revoir en un seul morceau.

Enfin, en quatre.

Obélik récupéra l'embarcation et rejoignit les autres dehors. Azaka accueillit leur investissement avec satisfaction : le canoë faciliterait considérablement le trajet. Elle prit la tête du groupe et les conduisit hors de Port Nyanzaru, vers les plaines qui séparaient encore la cité de la véritable jungle.

Au passage, ils croisèrent plusieurs Tabaxi, des êtres félins dont l'apparence rappela aussitôt à certains le prospectus de River Mist et Flask of Wine. Les aventuriers ne s'attardèrent pas. Au-delà des dernières habitations, les bruits de la ville s'effacèrent et l'immensité de Chult s'ouvrit devant eux.

Ils appliquèrent le répulsif acheté la veille avant que les insectes ne deviennent un problème. La lotion grasse ne rendait personne particulièrement élégant, mais elle devait les protéger pendant toute la journée. Azaka les conduisit jusqu'à la rivière et leur indiqua qu'ils devraient progresser vers le sud pendant plusieurs jours.

La guide les avertit que les ramifications du cours d'eau pouvaient facilement tromper un voyageur. Sous la canopée, la lumière traversait encore les feuilles, mais le soleil ne permettrait plus toujours de distinguer les points cardinaux. Même si elle connaissait le chemin, chacun devrait rester attentif aux embranchements et aux repères.

Ils mirent le canoë à l'eau. Obélik s'installa seul sur le banc central afin de répartir son poids, tandis qu'Azaka prit place à l'arrière pour diriger l'embarcation. Les autres se répartirent entre l'avant et le dernier banc disponible, puis les pagaies plongèrent dans le courant.

Port Nyanzaru s'éloigna derrière eux.

\illustration[0.9\textwidth]{2026-10-06/18220b70-2325-428d-9c61-5df12a5d47a3.png}
{\textit{Guidés par Azaka, les aventuriers quittèrent Port Nyanzaru.}}
{depart_port_nyanzaru}

\scenebreak

La plaine céda progressivement la place à une végétation plus épaisse. Les arbres se rapprochèrent des rives, leurs branches s'entremêlèrent au-dessus de l'eau et le fleuve commença à serpenter sous une voûte verte. Autour du canoë, chaque bruissement semblait pouvoir dissimuler une créature inconnue.

Après un virage, Kael perçut un bruit différent du murmure de la rivière et des cris d'animaux. Quelque chose glissait et frottait au-dessus de l'eau. En observant plus attentivement le passage devant eux, il distingua de minces filaments qui accrochaient par instants la lumière.

Une toile presque invisible barrait la rivière d'une rive à l'autre.

Kael prévint immédiatement ses compagnons. S'ils entendaient les créatures tisser depuis le canoë, elles ne devaient pas être petites. Azaka confirma son inquiétude et leur ordonna de se préparer.

Deux araignées gigantesques apparurent parmi les arbres.

Obélik réagit le premier. Il saisit l'une de ses haches à main et la lança vers la créature la plus proche, mais l'arme passa à côté de sa cible et disparut dans l'eau. Le goliath grogna, sauta hors du canoë et nagea jusqu'à la rive, découvrant trop tard que le fond se trouvait bien plus bas qu'il ne l'avait espéré.

Kael désigna l'une des araignées comme sa proie, puis banda son arc. Sa première flèche traversa l'espace entre les branches et s'enfonça profondément dans le corps velu de la créature. Azaka bondit à son tour hors de l'embarcation pour gagner la berge et se rapprocher des monstres.

Destos invoqua la lumière afin de bénir plusieurs de ses compagnons. Une énergie rassurante les enveloppa tandis qu'ils se dispersaient autour du canoë. Morrÿs lança ensuite son \textit{Fire Bolt}, mais la petite motte de feu passa entre les longues pattes de l'araignée et se perdit dans la végétation.

\illustration[0.9\textwidth]{2026-10-06/b8ef48aa-7f5f-45c4-9dd6-923d2b4a251a.png}
{\textit{Le groupe affronta deux araignées géantes au milieu de leur toile.}}
{combat_araignees_geantes}

La riposte ne tarda pas. L'une des créatures projeta un paquet de toile vers Kael. L'orc se baissa juste à temps et sentit les filaments siffler près de son visage avant de se coller derrière lui. L'autre araignée se rua sur Azaka et planta ses crochets dans son bras.

Kael vit la morsure entamer profondément sa chair.

Il ne vit pas une seule goutte de sang.

La plaie commençait déjà à se refermer lorsque l'araignée retira ses crochets. Azaka ne donna aucune explication et poursuivit le combat comme si cette blessure n'avait jamais existé. Au milieu de l'affrontement, Kael n'eut pas le loisir de lui poser la moindre question.

Obélik atteignit enfin la première créature. Sa grande hache s'abattit avec une violence qui arracha un cri strident au monstre. Depuis le canoë, Kael tenta une seconde flèche, mais l'embarcation tanguait sous les déplacements de ses compagnons et son tir se perdit dans les feuillages.

\begin{quote}\itshape

--- Arrêtez de bouger derrière !

\end{quote}

Azaka dégaina son Yklwa et se jeta dans la mêlée. Son premier coup manqua sa cible, mais elle corrigea aussitôt son mouvement. Lorsque l'araignée abaissa son corps pour frapper, la guide trancha l'une de ses pattes et la força à reculer.

Destos appela une flamme sacrée sur la créature blessée. L'araignée esquiva la lumière au dernier instant. Morrÿs eut davantage de réussite : son rayon glacé frappa le monstre et ralentit ses mouvements, couvrant une partie de son corps d'une morsure de givre.

L'araignée tenta de se jeter sur le canoë, mais son bond échoua. Obélik profita de l'ouverture et abattit une nouvelle fois sa hache sur elle. Malgré les blessures qui s'accumulaient, la créature refusait encore de s'effondrer.

L'araignée s'étant suffisamment rapprochée, Kael abandonna son arc. Il tira son épée courte et son cimeterre. La première lame pénétra entre deux plaques de chitine ; la seconde traça une profonde entaille dans l'abdomen déjà marqué par sa flèche.

Azaka rejoignit le combat d'un bond. Une patte griffue tenta de la happer lorsqu'elle se déplaça, sans parvenir à la toucher. Deux frappes rapides de son Yklwa achevèrent enfin l'araignée, qui bascula dans la rivière et disparut sous la surface.

La seconde créature restait encore debout. Destos tenta de nouveau de l'atteindre avec sa lumière, mais elle se déroba. Le \textit{Fire Bolt} de Morrÿs passa lui aussi à côté d'elle.

L'araignée choisit Obélik.

Ses crochets se refermèrent sur le goliath. Il contracta ses muscles et sa peau sembla prendre la dureté de la pierre, absorbant l'essentiel de la morsure. Le venin pénétra tout de même dans ses veines, lui arrachant une grimace, mais il résista suffisamment pour rester debout.

Obélik leva sa grande hache.

La lame s'abattit sur la tête de l'araignée et la sépara net de son corps. Les longues pattes continuèrent de se contracter pendant quelques instants, animées de soubresauts désordonnés, avant de s'immobiliser enfin sur la berge.

Le silence revint progressivement sur la rivière. Les fils tendus entre les arbres remuaient encore sous la brise, mais plus rien ne s'avançait pour défendre la toile. Le groupe venait de remporter son premier combat depuis son entrée dans la jungle.

Ils n'avaient quitté Port Nyanzaru que depuis quelques heures.
